% !TeX root = t3_3000.tex
        
\section{Effiziente Objekterkennung zur Analyse am Dart-Gerät}

Um eine automatisierte Prüfung des Wurfabstands sowie eine Unterscheidung zwischen geworfenen und gesteckten Pfeilen umzusetzen, ist eine ressourcenschonende Objekterkennung notwendig. 
Durch die Lokalisierung des Spielers im Kamerabild kann mathematisch geprüft werden, ob die Distanz zum Dart-Gerät in etwa den Regeln entspricht. Um den Anforderungen an eine geringe CPU- und Arbeitsspeicherauslastung auf der Zielhardware gerecht zu werden, 
müssen die untersuchten Verfahren eine niedrige algorithmische Komplexität aufweisen \parencite{_14, _2025_amd}.\newline

Techniken wie \enquote{Support Vector Machines} (SVM) in zusammenarbeit mit \enquote{Histogram of Oriented Gradients} (HOG) sind dabei vorteilhaft, da sie die CPU weniger belasten als tiefe neuronale Netze \parencite{_14}. Dies stellt sicher, dass die Hardware, 
genügend Kapazitäten für den zeitgleichen Spielbetrieb und die Videoverarbeitung behält \parencite{_2025_amd}.
		